\chapter{Two circulars, from business meaning to executable controls}
\label{ch:circular}


The two cases answer different questions. Circular 23EC35 shows a complete
selected transaction profile: qualification, evidence, consent, execution and
change. Circular 23EC46 then exposes why successful execution cannot settle an
uncertain classification. Both examples use retained public editions and
synthetic facts. Their different approval states are part of the lesson.

% BEGIN SOURCE UNIT P-01-problem:02
\section{SPI qualification: the conditions a wealth flag conceals}
\label{merge:P-01-problem:02}

\label{sec:spi}

The joint circular of 28 July 2023 supplies an unusually useful test case because
its apparent simplicity conceals several different kinds of rule. Annex 1
defines the qualifying criteria, eligible investment transactions, information
and consent requirements, and continuing controls. Annex 2 clarifies difficult
implementation cases. The examples below are synthetic and illustrate the cited
provisions; they do not complete an SPI assessment or authorize a transaction
\citep{sfcspi,sfcspiannex1,sfcspiannex2}.

\caseheading{A one-word error changes the financial condition}
Annex 1, paragraph 3.1 provides two alternative financial routes: a portfolio of
at least HK\$40 million, or net assets excluding the primary residence of at least
HK\$80 million, including the stated foreign-currency equivalents. Consider a
client with an eligible portfolio valued at HK\$35 million and qualifying net
assets of HK\$90 million. The second route meets this financial condition.
Replacing the alternative with a conjunction rejects the client incorrectly.
Replacing either inclusive comparison with a strict comparison rejects a client
exactly at that threshold. Both defects are easy to miss when test data use
wealth amounts far from the boundary \citep[Annex 1, para.~3.1]{sfcspiannex1}.

The numbers also require definitions. The value of assets held at the bank is
not automatically the qualifying portfolio, and total wealth is not automatically
net assets excluding the primary residence. Paragraph 3.3 specifies the net-asset
calculation. In a simple synthetic case with HK\$120 million of assets, including
a HK\$30 million primary residence, and HK\$15 million of liabilities, that
calculation gives HK\$75 million: $120-15-30$. Omitting the residence gives
HK\$105 million and changes the result of the net-asset route. The portfolio
route must still be assessed independently. The bank must also establish how the
source definitions map to its valuation and ownership records
\citep[Annex 1, paras.~3.2--3.4]{sfcspiannex1}.

\caseheading{Ownership and missing evidence can reverse the result}
For a joint portfolio with persons other than the client's associate, paragraph
3.2(c) refers to the client's share specified in a written agreement, or an equal
share in the absence of that agreement. Suppose a two-holder portfolio is worth
HK\$60 million and the written allocation gives the client 20 percent. The
relevant share for this illustration is HK\$12 million, not the whole account
and not an automatically assumed half. The rule for a joint account with an
associate is separately stated. A generic ``joint-account multiplier'' loses
this distinction \citep[Annex 1, para.~3.2]{sfcspiannex1}.

An unavailable document is also different from an established absence of an
agreement. The bank must decide what evidence establishes each fact. If a feed
returns a null portfolio value, replacing null with zero creates a substantive
conclusion from a data failure. If qualifying net assets are known to be
HK\$90 million, the financial condition can still be established by that route.
If they are HK\$70 million, the unknown portfolio value can change the answer;
the financial condition remains unresolved. The product must express that
distinction directly.

\caseheading{Wealth does not establish suitability for streamlined treatment}
Paragraph 1 combines professional-investor status, the financial criteria,
knowledge or experience, and investment objectives at the relevant date.
Paragraph 4.1 requires reasonable satisfaction concerning sophistication and
sets out alternative supporting criteria. Paragraph 5.1 excludes conservative
clients from SPI treatment for this guidance. A wealthy client whose objective
is capital preservation therefore demonstrates why a field named
\texttt{wealth\_qualified} must not automatically populate
\texttt{streamlined\_approach\_available}
\citep[Annex 1, paras.~1, 4--5]{sfcspiannex1}.

Some criteria are countable; others require assessment. A degree, a qualification,
relevant work experience and prior transactions are distinct routes. A translator
should not invent a score that trades off wealth against sophistication or
investment objectives. The assessment and its supporting evidence need an owner,
date and review record. Paragraph 2 also supplies a separate route for a qualifying
investment-holding corporation wholly owned by SPIs. It cannot be represented by
treating every client identifier as an individual.

\caseheading{Experience and consent belong to a product category}
Suppose the client relies on the transaction-experience route and has completed
five bond transactions in the relevant three-year period. That history does not
automatically establish the same route for an unrelated product category.
Paragraph 7.3 connects category eligibility to the alternative qualification and
experience routes; paragraph 7.2 requires the client's choice of categories to be
documented and category information to be provided. A global client flag can
therefore remain true while a proposed category is outside the selected scope
\citep[Annex 1, paras.~4.1, 7.1--7.3]{sfcspiannex1}.

The identity of the person matters too. Annex 2, FAQ 2 explains the assessment
where a client acts through a power of attorney. The attorney's sophistication
cannot simply replace the client's qualifying attributes. In software terms,
the customer, account holder, representative, beneficial owner and person giving
an instruction are different roles. Collapsing them into one record can make
every subsequent logical step internally consistent and still apply it to the
wrong person \citep[Annex 2, FAQ 2]{sfcspiannex2}.

\caseheading{Exposure monitoring is an event rule}
Paragraph 8.3 permits two monitoring approaches. One concerns gross exposure on
execution of transactions under the streamlined approach. The other uses
designated accounts and checks the specified condition after a top-up or deposit
of new funds. Annex 2, FAQ 6 explains how these approaches operate, including
transferred exposures and restrictions on additional funds when a designated
account exceeds the threshold. FAQ 7 addresses leverage
\citep[Annex 1, para.~8.3; Annex 2, FAQs 6--7]{sfcspiannex1,sfcspiannex2}.

A single rule that rejects every transaction whenever an account's current value
exceeds a threshold loses the distinction between these approaches. In the
designated-account example, the FAQ allows continued execution above the
threshold subject to its conditions restricting new funds. The monitoring mode,
movement of funds, gross exposure and prior events all matter. An appreciated
position, a leveraged purchase, a cash withdrawal and a transferred non-cash
position must not be reduced to the same change in account market value. This is
a strong reason to model event histories as well as decision tables.

\caseheading{Consent and continuing obligations outlive the initial decision}
Paragraph 13 calls for prior written agreement and an explanation of the ability
to withdraw from, or change, the arrangement. Paragraph 14 requires annual review
and reminders. Assessment and agreement are therefore not permanent facts. A
transaction considered before receipt of a withdrawal instruction and one
considered afterwards require different state snapshots. A review reminder being
queued is different from being sent; being sent is different from the review
being completed \citep[Annex 1, paras.~12--14]{sfcspiannex1}.

The streamlined approach also retains specific information and conduct
requirements. The solicited or recommended transaction provisions and the
unsolicited complex-product provisions have different contexts. A broad
``skip suitability and disclosure'' switch would conceal these remaining duties
\citep[Annex 1, paras.~9--11]{sfcspiannex1}. The prototype should return the
relevant decision and the associated tasks, explanations and evidence requests,
rather than a single apparently comprehensive compliance flag.
% END SOURCE UNIT P-01-problem:02

% BEGIN SOURCE UNIT P-01b-complete-dry-run:00
\section{A complete worked change: circular 23EC35 to a Java decision}

The retained SPI example provides the first end-to-end execution lesson.
Its original standalone compiler is called SPI-Demo1. The later local workbench
extends that lesson through authenticated synthetic review, attributed events,
release and archive restoration. This chapter first makes the business meaning
and the small executable example explicit; Chapters~\ref{ch:verification} and
\ref{ch:implementation} distinguish the two interfaces and their actual evidence.


\label{merge:P-01b-complete-dry-run:00}

\label{sec:dry-run}

The running source is the joint SFC--HKMA circular of 28 July 2023,
\emph{Streamlined approach for compliance with suitability obligations when
dealing with sophisticated professional investors}. Its main text introduces
the approach; Annex 1 sets out the guidance; Annex 2 answers implementation
questions. This example follows those documents through collection,
interpretation, evidence mapping, evaluation, Java generation and an operational
change. The facts about the client are invented for the demonstration. The
regulatory provisions are identified individually so that those two sources of
information remain distinguishable \citep{sfcspi,sfcspiannex1,sfcspiannex2}.
% END SOURCE UNIT P-01b-complete-dry-run:00

% BEGIN SOURCE UNIT P-01b-complete-dry-run:01
\subsection{The business question and the facts the bank would need}
\label{merge:P-01b-complete-dry-run:01}


At 10:00 Hong Kong time on 21 September 2026, an adviser proposes a solicited
transaction in an ordinary corporate-bond category for Ms Lee, a synthetic
individual client. The question is whether the bank may apply the selected
streamlining treatment to this transaction, subject to its other controls.
The example uses paragraph 8.3(a), which monitors gross exposure at execution.
It does not use the designated-account alternative in paragraph 8.3(b).

Assume the bank has already assessed Ms Lee as an individual professional
investor under the applicable imported definitions. Her qualifying portfolio is
HK\$40 million: HK\$35 million on her own account plus a documented 25 percent
share of a HK\$20 million joint portfolio with a non-associate. Her total assets
are HK\$120 million, liabilities HK\$15 million and primary residence HK\$30
million, giving HK\$75 million of net assets excluding the residence. The
portfolio route meets paragraph 3.1 exactly; the net-asset route does not.
The ownership agreement, valuations and assessment would be separate evidence
records in the bank. The example's normalized inputs already incorporate them.

Ms Lee has no relevant degree, professional qualification or qualifying year of
work experience in this example. She instead has five transactions in the
chosen category, dated 8 January and 12 June 2024, 14 February and 7 November
2025, and 20 August 2026. These dates lie comfortably inside the relevant
three-year window; the example does not settle a disputed endpoint convention.
The bank has separately recorded its reasonable satisfaction about her
sophistication and assessed her objectives as non-conservative. Five transactions
are evidence for a specified route, not an automatic substitute for that
assessment \citep[Annex 1, paras.~4.1, 5.1, 7.3]{sfcspiannex1}.

The client has selected this category, received the category information and
required explanations, and given the prior written acknowledgment for the
arrangement. The agreed streamlining threshold is HK\$10 million. Existing
gross exposure is HK\$8 million; the proposed transaction adds HK\$2 million,
including leverage, bringing the monitored exposure to HK\$10 million. The
control compares that gross amount with the threshold, not the cash margin
posted for the purchase. The threshold rationale and assessment are retained;
the example treats the annual review and reminders as current
\citep[Annex 1, paras.~6--8, 12--14; Annex 2, FAQ~7]{sfcspiannex1,sfcspiannex2}.

Finally, current offering documents have been delivered. There is no unanswered
request for product explanation or material query. The transaction has been
assessed for the outsize or material-transaction warning requirement; this
synthetic case does not require that warning. These are retained controls.
They do not disappear because the wealth condition passed
\citep[Annex 1, paras.~8.4, 10.3]{sfcspiannex1}.

The facts are intentionally close to two boundaries. A one-cent reduction in
qualifying portfolio, with the other facts fixed, defeats the financial
condition. A one-cent increase in projected gross exposure defeats the selected
execution-monitoring condition. A test suite containing only a very wealthy
client buying a very small position would miss both defects.
% END SOURCE UNIT P-01b-complete-dry-run:01

% BEGIN SOURCE UNIT P-01b-complete-dry-run:02
\subsection{Collect the whole instruction before extracting rules}
\label{merge:P-01b-complete-dry-run:02}


The intake record contains reference number 23EC35, the public document URL,
issuer, publication date, language, retrieval time and content digest. It
links the main text to both annexes and records their digests separately.
The annex PDFs are retained as immutable bytes. Extracted text is a derivative,
with page breaks, its own digest and exact character offsets. A source anchor
therefore identifies a particular passage in a particular preserved edition,
even if the public web page later changes.

For example, the financial condition is anchored to Annex 1 paragraph 3.1.
The supplied source record includes the PDF digest beginning
\texttt{f83719981d74}, the derivative digest, the page and the span digest.
The validator checks those against the actual retained files. Changing one
character in the extraction invalidates an old offset-based anchor until the
passage is matched and reviewed again. A language model cannot repair an anchor
by supplying a plausible-looking quotation.

Collection also produces a dependency list. Annex 1 imports definitions from the
Professional Investor Rules and the Code of Conduct; the opening footnote
identifies suitability requirements in Code paragraphs 5.2 and 5.5(a). Annex 2
and Annex 1's footnotes refer to further product guidance. The demonstration
accepts reviewed input assessments for these imported matters. A production
release must name the adopted source editions, affected legal entity and
business activity. The synthetic bundle's September 2026 validity interval is
a test setting, not a finding about the circular's legal commencement or the
continued applicability of every imported provision.

An agent should next produce a provision inventory, before proposing a formula.
The following inventory covers the numbered guidance and FAQ families. A row
can lead to a decision rule, an evidence calculation, a human assessment or a
continuing task. ``Outside this profile'' means a separate implementation is
needed; it never means the passage was irrelevant and discarded.

\begingroup\small
\begin{longtable}{P{0.14\textwidth}P{0.37\textwidth}P{0.39\textwidth}}
\caption{Disposition of the complete circular for the worked profile. Subparagraphs
within each stated range remain part of that row.}\label{tab:circular-inventory}\\
\toprule Source & Substantive instruction & Treatment in the worked change \\\midrule
\endfirsthead
\toprule Source & Substantive instruction & Treatment in the worked change \\\midrule
\endhead
Main text; Annex 1 opening and footnote 1 & Introduce streamlining, retain conduct and risk controls, and import definitions and suitability requirements. & Record authority, perimeter and dependencies. The Java result covers selected conditions only; the bank retains its wider control chain.\\
1--2 & Individual criteria at the relevant date; separate investment-holding corporation route. & Individual PI assessment is an input. A corporate client returns outside profile; ownership and qualifying-owner aggregation require a separate rule family.\\
3.1(a)--(b) & Inclusive alternative wealth thresholds and currency equivalents. & Execute exact HK-cent comparisons joined by OR. Foreign currency requires an adopted conversion mapping before evaluation.\\
3.2(a)--(d), 3.3, 3.4(a)(i)--(iii), (b) & Portfolio and asset/liability attribution; residence exclusion; written or equal joint shares under the stated circumstances. & Reviewed normalization produces portfolio and net assets. A missing agreement is unresolved evidence, not established absence. Preserve each ownership calculation.\\
4.1(a)--(d), 5.1 & Reasonable sophistication assessment, alternative supporting routes and non-conservative objectives. & Execute the alternative routes and require the dated assessment and objectives. Do not substitute the representative's attributes.\\
6, 7.1--7.3 & Restrict transactions to selected categories and threshold; define categories, provide information, and apply category qualification. & Require selected category, resolved category explanations and information. Count same-category transactions or establish another qualification route.\\
8.1--8.2 & Client specifies a threshold appropriate to circumstances; bank retains the setting and rationale. & Accept an agreed, reviewed absolute HKD threshold and rationale. A percentage-of-AUM arrangement needs a versioned conversion to the same input.\\
8.3(a) & Monitor gross exposure on execution. & Compare projected exposure including the candidate transaction with the agreed limit; host must reserve exposure atomically before using the decision.\\
8.3(b) & Alternative monitoring for designated accounts and new funds. & Explicitly outside profile. Implement its own account, transferred-position and top-up events before enabling it.\\
8.4--8.5 & Detect outsize or material transactions and warn; review threshold compliance annually and alert if exceeded. & Require transaction assessment and conditional warning. Schedule review/alert tasks; a reviewed composite reports annual controls current.\\
9; 10.1--10.4 & Define solicited-transaction relief, with explanation exceptions and offering documents retained. & Execute the solicited profile. The result identifies applicable relief; it does not disable every suitability or disclosure control.\\
11.1--11.5; footnotes 2--4 & Different relief for unsolicited complex-product transactions, including document alternatives and warning arrangements. & Outside profile. Preserve the distinct document/evidence branches, bond information alternatives and annual-warning requirement for a later implementation.\\
12.1--12.3 & Reasonable satisfaction, written assessment and correspondence recording category and threshold choices. & Dated assessments and document references feed named facts; records remain independently inspectable.\\
13.1(a)--(d), 13.2 & Prior written agreement and explanation of consequences, categories, threshold and withdrawal/amendment rights. & A checklist-backed acknowledgment input plus a replayed consent state. Every item in the checklist must be evidenced; a signature alone is insufficient.\\
14.1--14.2 & Annual continued-qualification/agreement review and written reminders covering consequences, categories, threshold, breaches and rights. & Ongoing tasks and reviewed status. Suspending streamlining when review is overdue is a proposed additional bank restriction.\\
Annex 2, FAQ 1 & Appropriate reliance on disclosures after reasonable KYC efforts; clarify inconsistency. & Define admissible evidence and escalation; do not invent a blanket requirement for independent third-party confirmation.\\
FAQ 2 & Client and power-of-attorney distinctions. & Attribute qualifications to the client. Apply the FAQ to transaction-history evidence without copying the attorney's sophistication.\\
FAQs 3--4, including notes & Category design, difficult product distinctions and category information. & Bank-owned category taxonomy, versioned mappings and information checklist; uncertain mapping requests review.\\
FAQ 5 & Threshold should reflect the client's circumstances; bank-held AUM may be a limited part of wealth. & Recorded client circumstances and threshold rationale. No universal percentage is invented.\\
FAQ 6 & Explain both monitoring approaches, including designated-account transfers, appreciation and new funds. & Use execution monitoring here. Keep alternative-account rules visible; do not infer forced unwinding merely from exceeding its threshold.\\
FAQ 7 & Gross exposure includes leverage. & Feed leverage-inclusive projected exposure to the execution comparison.\\
\bottomrule
\end{longtable}
\endgroup

This inventory is the first useful review deliverable. It exposes omissions
before a test generator starts exploring the rules that happened to be encoded.
The main text, annexes and incorporated definitions are the review boundary;
the numeric threshold is only one part of it.
% END SOURCE UNIT P-01b-complete-dry-run:02

% BEGIN SOURCE UNIT P-01b-complete-dry-run:03
\subsection{Record the interpretation and who supplies each fact}
\label{merge:P-01b-complete-dry-run:03}


The agent's draft should separate source statements from proposed bank choices.
For paragraph 3.1 it records an inclusive alternative and links both comparisons
to the paragraph. For paragraph 13 it records an arrangement-specific written
acknowledgment and the withdrawal right. For paragraph 14 it records the annual
duties. It then proposes, separately, that the bank suspend use of the streamlined
route when its annual-review status is not current. A compliance reviewer can
accept, amend or reject that last proposal without losing the source duty.

Two unresolved issues are deliberately visible. First, the bank must settle
ownership, currency and valuation mappings for wealth and exposure. Second,
the calendar conventions for transaction windows and annual deadlines must be
adopted. A complete source-based interpretation may still rely on human
assessment: ``reasonably satisfied'' and ``material transaction'' do not acquire
objective numerical definitions simply because the system needs an input.

In the production workflow, the preparer submits a specific bundle digest and
supporting examples. An authorized reviewer other than the preparer accepts that
exact version, with the perimeter and unresolved issues visible. Editing a rule
changes the digest and makes the previous acceptance insufficient for release.
The delivered example carries \texttt{DEMONSTRATION\_ONLY} throughout: its
synthetic assessments and proposed bank choices have no invented reviewer
approval. Draft evaluation is useful before approval; production eligibility
requires the release controls described later.

\begingroup\small
\begin{longtable}{P{0.30\textwidth}P{0.60\textwidth}}
\caption{Input contracts for the demonstrated Java profile. A known value always
has evidence, a validity interval and a recording time.}\label{tab:demo-inputs}\\
\toprule Input or group & Meaning and intended owner \\\midrule
\endfirsthead
\toprule Input or group & Meaning and intended owner \\\midrule
\endhead
Profile flags & \texttt{profile\_individual}, \texttt{profile\_solicited}, \texttt{profile\_execution\_monitoring}: the host's classified client, activity and monitoring mode. False selects an unsupported profile; unknown requests review.\\
\texttt{individual\_pi} & Compliance assessment under the adopted Professional Investor Rules and Code definitions, for this legal client and date. It is distinct from merely being an individual.\\
\texttt{portfolio}; \texttt{net\_assets\_ex\_home} & Valuation/operations supplies exact HK cents after reviewed ownership, liabilities, residence and currency treatment. Preserve the calculation inputs and source evidence.\\
Qualification fields & \texttt{relevant\_degree}, \texttt{professional\_qualification}, \texttt{relevant\_year\_of\_work}; alternatively \texttt{category\_transactions\_3y}. Client-level evidence is mapped to the selected category and relevant period.\\
Assessment fields & \texttt{reasonable\_sophistication\_assessment}, \texttt{non\_conservative\_objectives}, \texttt{written\_assessment\_retained}. Compliance/relationship management supplies the assessment and records; missing review is unknown.\\
Category fields & \texttt{category\_selected}, \texttt{category\_information\_delivered}, \texttt{category\_explanation\_resolved}. Client choice and delivery evidence must match the transaction's versioned category.\\
Acknowledgment and consent & \texttt{prior\_written\_acknowledgment\_complete} covers every required item in 13.1--13.2; \texttt{active\_consent} comes from the client/category consent history. Document completeness and present consent are different facts.\\
Threshold and exposure & \texttt{streamlining\_threshold}, \texttt{projected\_gross\_exposure} are HK cents; \texttt{threshold\_rationale\_retained} is an assessment/record check. Order management supplies a consistent exposure snapshot including leverage.\\
\texttt{annual\_review\_current} & Reviewed composite covering continued qualification/agreement, threshold review, required reminders and breach alerts. The demo consumes this status; it does not calculate annual deadlines or send reminders.\\
Product information & \texttt{current\_offering\_documents\_delivered}, \texttt{explanation\_requested\_or\_material\_query}, \texttt{explanation\_delivered}. Delivery/interaction records determine whether the remaining request has been answered.\\
Transaction warnings & \texttt{transaction\_materiality\_assessed}, \texttt{outsize\_or\_material\_transaction}, \texttt{material\_warning\_delivered}. The assessment is required; a warning is required if the assessment says so.\\
\bottomrule
\end{longtable}
\endgroup

Composite inputs are integration boundaries, not permission to hide missing
work. For example, the acknowledgment evidence record must link the prior
written agreement, qualification explanation, chosen categories and threshold,
consequences and withdrawal/amendment explanations. A rejected or incomplete
checklist makes the composite false or unknown under the bank's adopted evidence
policy. The reviewer can open the constituent items. The generated decision
library does not authenticate documents or decide whether an explanation was
adequate.
% END SOURCE UNIT P-01b-complete-dry-run:03

% BEGIN SOURCE UNIT P-01b-complete-dry-run:04
\subsection{Construct and evaluate the executable specification}
\label{merge:P-01b-complete-dry-run:04}


The worked bundle has ten named Boolean rules. The financial rule performs the
two monetary comparisons. The category-qualification rule joins the four
supporting routes. The assessment rule combines individual PI status, finance,
category qualification, reasonable satisfaction, objectives and written records.
A separate review-policy rule records the proposed bank restriction under its
own bank-policy interpretation. The arrangement rule checks category choice and information, acknowledgment,
active consent, threshold rationale and the proposed annual-review restriction.
The threshold rule compares limit with projected gross exposure. The explanation
and warning rules express conditional retained actions. The retained-information
rule combines these with delivery of current offering documents. The root rule
requires assessment, arrangement, threshold and retained information within the
selected profile.

Every reference points to a named, unconditional subrule; the root alone selects
the profile. The graph is acyclic, so evaluation terminates after a finite number
of expressions. Money and counts are exact integers. Unknown evidence is
propagated by the specified three-valued operations; contradictory supplied
records return a conflict before evaluation. The compiler type-checks the tree
and rejects unsupported operators instead of silently choosing a different
meaning.

For Ms Lee, the trace reads as follows. The portfolio comparison is true at
HK\$40 million; the net-asset comparison is false at HK\$75 million; finance
is therefore true. Five category transactions establish the qualification route.
The other reviewed assessment and arrangement conditions are true. HK\$10
million is at most the agreed HK\$10 million, so threshold monitoring succeeds.
The explanation request is false, the transaction is assessed as not requiring
the material-transaction warning, and offering documents are delivered. The root
returns \texttt{STREAMLINING\_CONDITIONS\_MET}. The result includes the rule
trace, source-anchor identifiers, rule-bundle and fact-snapshot digests, both
evaluation timestamps and a digest of the result itself.

This successful conclusion remains narrow. It identifies the selected
streamlining conditions under the supplied assessments and proposed review
restriction. The host must still apply other relevant product, credit, market,
financial-crime, conduct and order-execution controls. Its own release process
decides which approved control version is usable. A demonstration result is
never automatically promoted by setting a Boolean input named ``approved''.

\begingroup\small
\begin{longtable}{P{0.38\textwidth}P{0.19\textwidth}P{0.33\textwidth}}
\caption{Distinguishing cases derived from the running interpretation. Other
inputs remain as in the successful synthetic case unless stated.}\label{tab:dry-cases}\\
\toprule Changed fact & Result & Why it changes the operational response \\\midrule
\endfirsthead
\toprule Changed fact & Result & Why it changes the operational response \\\midrule
\endhead
Portfolio HK\$39,999,999.99; net assets HK\$75m & Standard process & Both financial alternatives are refuted.\\
Portfolio HK\$35m; net assets HK\$90m & Conditions met & The alternative net-asset route suffices.\\
Portfolio missing; net assets HK\$90m & Conditions met & Missing portfolio is reported, but does not block the established alternative.\\
Portfolio missing; net assets HK\$70m & Review required & Portfolio evidence can change the answer. It is a blocking input.\\
Conservative objectives & Standard process & Wealth does not overcome paragraph 5.1.\\
Four category transactions; no other qualifying route & Standard process & The countable route is not established.\\
Chosen category does not match & Standard process & General SPI attributes do not extend the selected categories.\\
Gross exposure HK\$10,000,000.01 & Standard process & The execution-monitoring condition is exceeded.\\
Written acknowledgment incomplete, or consent withdrawn & Standard process & Prior documentation and current consent are both required.\\
Client requests explanation; no answer recorded & Standard process & Paragraph 10.3 retains the requested explanation.\\
Material transaction; no warning delivered & Standard process & The applicable paragraph 8.4 warning remains outstanding.\\
Two inconsistent portfolio records & Review required & An evidence conflict is not resolved by choosing the more convenient amount.\\
Designated-account monitoring selected & Outside profile & Its different paragraph 8.3(b)/FAQ 6 logic has not been executed by this library.\\
\bottomrule
\end{longtable}
\endgroup

The associated actions are equally concrete. A missing portfolio requests the
valuation and ownership evidence. An unanswered explanation routes a task to the
adviser and retains the interaction record. An exceeded execution threshold
routes the order to the standard process or a separately reviewed change in
arrangement. An evidence conflict routes to its data owner. These are proposals
for the host workflow, not a generated instruction to place or cancel a trade.
% END SOURCE UNIT P-01b-complete-dry-run:04

% BEGIN SOURCE UNIT P-01b-complete-dry-run:05
\subsection{Generate, build and call the Java release candidate}
\label{merge:P-01b-complete-dry-run:05}


The demonstrator stores the typed tree in \path{spec/spi-control.bundle.json}
under \path{examples/java-dry-run/}. A deterministic
compiler walks each node and emits a call to its corresponding Java operation.
An integer literal becomes a \texttt{BigInteger}; a fact reference becomes a
typed fact lookup; \texttt{any} becomes the explicitly three-valued alternative;
a rule reference becomes a memoized Java method. Source identifiers and the
bundle digest are compiled into the class. There is no language-model call,
network retrieval or rule-file parsing during a transaction decision.

The generated class and its small support library form the JAR
\path{spi-controls-demo.jar}. A separate host example is
then compiled against that JAR. This tests the actual deployment boundary:
the caller sees the public library API and receives a structured decision.
Section~\ref{sec:java} specifies the API and how a bank adapter must prepare
its inputs. The same build compares complete Java decision results with the
independent Python reference, rather than merely checking that compilation
succeeded.

Three deliberate generated-code mutations make the test useful. Changing the
portfolio comparison to strict rejects the inclusive boundary. Replacing the
wealth alternative with a conjunction rejects the alternative route. Bypassing
the consent condition permits the withdrawal case. The build compiles and runs
each mutated rule class against the same client harness and requires at least
one independently specified status to fail for each mutation. Matching hashes
alone are not considered detection of a semantic defect.
% END SOURCE UNIT P-01b-complete-dry-run:05

% BEGIN SOURCE UNIT P-01b-complete-dry-run:06
\subsection{Withdraw consent, replay the decision, then change the rule}
\label{merge:P-01b-complete-dry-run:06}


At 10:05 Hong Kong time, suppose Ms Lee withdraws. The event is recorded at
10:06. A replay has two timestamps: the time being assessed and the cutoff for
what the system knew. The 10:00 decision remains reproducible with its original
facts and knowledge cutoff. A decision at 10:10 using a complete history known
at 10:10 sees the withdrawal and returns the standard-process result.
The Java host demonstration performs this sequence by replaying the consent
events, replacing the consent fact with the evidenced inactive value, and
evaluating the same generated control again.

If the event feed is incomplete, absence of a withdrawal record cannot establish
consent. The replay interface therefore requires a completeness assertion for
the relevant history window and the time that assertion became known. A record cannot certify completeness
through a future instant; the replay rejects that premise as unresolved. In the
late-arrival case, a history complete only through 10:04 does not support a
10:10 conclusion, even if the last visible event is a grant. Identical event
identifiers and payloads are idempotent duplicates; differing payloads under
the same identifier conflict. Events at the same occurrence time require an
authoritative sequence. The demo covers these behaviors explicitly; a production
adapter must establish the completeness and sequencing assertions, not invent
them.

Now consider a later amendment. As an engineering test only, suppose an approved
rule version changes a threshold or retained information condition. This is not
a claim that the SFC has issued that change. The collector creates a new source
revision; the workbench marks affected definitions, rules, examples and Java
classes. The reviewer sees the semantic difference and the clients whose
synthetic outcomes change. The compiler creates a new bundle digest, generated
source and JAR. Release approval applies to those exact versions. The old
bundle, JAR and evidence remain available for replay; the bank changes the
active release atomically at the agreed effective instant.

Rollback is a software-release action, not reversal of the legal calendar.
An operationally defective JAR may be replaced with another implementation of
the same approved legal version. Restoring an older legal rule after its
applicability has ended would require a separate, valid legal basis. Historical
replay selects the version that applied to the historical question, rather
than whichever JAR happens to be active today.

The dry run thus reaches an actual executable Java boundary while exposing
the work that remains: bank adjudication of interpretation, imported definitions,
data and calendar mappings, and production integration. The reviewer application and the broader
engineering modules were subsequently implemented in the local MVP;
Chapter~\ref{ch:implementation} records their evidence and the distinct ensemble
extension that remains to be built.
% END SOURCE UNIT P-01b-complete-dry-run:06

% BEGIN SOURCE UNIT P-01-problem:04
\section{Marketing rules require exceptions and judgment}
\label{merge:P-01-problem:04}


The October 2023 distributor circular examines additional returns and incentives,
redemption restrictions and marketing of fund-related services. Its treatment of
guaranteed-return arrangements depends on the arrangement and the relevant legal
classification. A potential structured-product classification should produce a
supported assessment task; a language model should not resolve it from a slogan
alone \citep[paras.~5--11]{sfcmarketing}.

The circular's discussion of gifts expressly distinguishes discounts of fees or
charges. Losing that exception makes an apparently conservative implementation
incorrectly broad. Its redemption discussion distinguishes administrative
procedures or cut-off times from reducing a fund's dealing frequency. A daily
dealing fund offered with weekly redemption presents a different issue from a
different daily processing cut-off. These are excellent paired test cases because
they differ in the fact that determines the treatment
\citep[paras.~9--14]{sfcmarketing}.

The advertising discussion addresses misleading impressions and comparisons with
deposits. Keyword detection can nominate material for review, but it cannot
establish the whole classification: meaning depends on the presentation and the
service described. The system should retain the advertisement, the relevant
source passages and the reviewer's conclusion. That is a useful result even when
there is no defensible fully automated predicate
\citep[paras.~15--17]{sfcmarketing}.
% END SOURCE UNIT P-01-problem:04

% BEGIN SOURCE UNIT M02:00
\section{The source packet and the question it supports}
\label{merge:M02:00}

\label{sec:circular-packet}
The running example begins with an identifiable public document: SFC circular
23EC46, issued on 24 October 2023, concerning additional returns and other
services or arrangements offered when marketing SFC-authorised funds
\citep{sfcmarketing}. The retained experiment selected its paragraph-10 gift
restriction. The circular as a whole discusses more than that restriction. A
reader needs to know both what was selected and what remains outside the control
before a result from the control can be interpreted.

The experiment preserved the official source response, an extracted text
derivative and anchors connecting the selected rule to its source passages.
On 22 September 2026, a fresh retrieval matched the retained raw bytes. This
check establishes that the retained response matched what the retrieval route
returned at that time. It does not establish the complete current legal position
or resolve the effective versions of incorporated sources. A hash identifies
bytes; it does not by itself establish their authority, completeness or legal
effect.

The selected assessment concerns one proposed incentive component, offered by a
distributor in promoting a fund within the declared profile. It is a question
about a specific restriction. It is not a whole-offer compliance verdict. The
distinction is visible in the result vocabulary. A true result means that the
selected prohibition trigger is established. A false result means that this
trigger is not established as true on the supplied, resolved facts under the
specified semantics. It supplies no general permission to market the offer.

The retained packet explicitly leaves the applicable Code paragraph and cited
FAQ unresolved. This matters twice. Those sources may settle the meaning or
scope of a term; they may also establish that the selected profile is narrower
than the underlying requirement. Their absence is therefore a release-blocking
dependency, not a bibliographic inconvenience. The example is useful precisely
because it continues to expose that absence after its executable tests pass.

The book uses this historical, selected packet to teach implementation and
assurance. New commercial offers, cashback definitions and mixed-component
scenarios introduced below are synthetic investigations. They are not reported
SFC rulings. The unresolved references are not silently filled using a language
model's recollection of a Code provision. Doing so would convert a visible gap
into an untraceable premise.
% END SOURCE UNIT M02:00

% BEGIN SOURCE UNIT M02:01
\section{Reading the surrounding circular before isolating the rule}
\label{merge:M02:01}

\label{sec:circular-context}
A translation that begins by copying the most convenient sentence can miss why
that sentence is there. The circular's opening establishes its distribution
context. The next passages describe observed practices, including additional
features, standing investment arrangements and marketing. Those observations
help identify the activity being discussed, but they are not automatically
individual executable prohibitions. A source inventory should retain their
function without manufacturing a rule for every paragraph.

The discussion then turns to additional returns and their possible regulatory
consequences. Some questions concern the structure of an arrangement and public
promotion. Others concern impressions created about the underlying fund. The
gift discussion appears within that broader context. Consequently, a finding
that one separate classification is absent does not clear the gift issue. The
gift control is not conditional on first finding that every other regulatory
concern applies.

Later passages address redemption arrangements, administrative cutoffs,
advertising and the description of services. Some concerns require event or
workflow evidence; others require qualitative assessment of communication.
Those mechanisms differ from the paragraph-10 Boolean control. An adequate
implementation of the whole circular would need to preserve those differences.
It cannot simply accumulate a larger expression around the gift predicate and
call the circular complete.

The existing disposition inventory has 18 numbered paragraphs and four
footnotes. Its assignments are agent-authored and unreviewed. They include
context, separate controls, qualitative review and unresolved dependencies.
Paragraph 10 is the only numbered passage with an executable control in this
exercise. A future independent reviewer should challenge both the assignments
and the completeness of the inventory. Merely approving the final formula would
leave the earlier decision about what to implement unexamined.

This provides an immediate requirement for the ensemble. One path must inspect
the source independently of the candidates and ask what obligations, exceptions
and dependencies need a disposition. Another path may develop the selected
rule. Comparing those paths creates a place for an omitted requirement to
appear. If both start from the same already-filtered list of paragraphs, the
apparent redundancy begins too late.
% END SOURCE UNIT M02:01

% BEGIN SOURCE UNIT M02:02
\section{Naming the six inputs precisely}
\label{merge:M02:02}

\label{sec:circular-inputs}
The executable example has two scope inputs and four body inputs. Their names
are convenient for software, but their definitions do the substantive work.
The first scope input records whether the assessed actor is a distributor in
the declared profile. The second records whether the activity concerns promotion
of an SFC-authorised fund within that profile. Both are reviewed classifications.
The prototype does not infer them from the actor's name or the presence of the
word ``fund'' in a document.

The first body input records whether the proposed component is classified as a
gift. The second and third record a promotional link to a specific product and
to a particular product type. A link to funds from a particular fund house
requires an explicit classification where the relevant example is used; the
engine does not assume that every mention of a fund house creates such a link.
The fourth body input records whether the entire assessed component falls
within the fee-or-charge-discount exception.

The word ``entire'' in that final definition is an engineering expression of
the proposed component-level reading. It prevents an assessed voucher from
inheriting the classification of a separate fee waiver. It does not prove that
the decomposition is legally correct. A reviewer must establish the relevant
unit of assessment from the source and context. The definition and the review
issue must travel together until that question is resolved.

For a compact mathematical account, let $A$ denote established actor scope and
$S$ the fund/activity scope. Let $G$ denote the gift classification, $P$ the
specific-product link, $T$ the product-type link and $D$ the discount exception.
For now suppose these are known Boolean values. The body condition is
\begin{equation}
 B=G\land(P\lor T)\land\neg D.
 \label{eq:gift-body}
\end{equation}
It asks three things together: is the component a gift, does either promotional
route hold, and is the exception absent? The two alternatives occur inside one
parenthesized condition. Replacing that disjunction by a test of $P$ alone
changes the rule for offers linked only to a product category.

Scope is assessed separately. When both scope inputs are true, the body answers
the selected question. When a scope input is known false, the limited profile
returns \texttt{OUT\_OF\_SCOPE}. When scope is unresolved, it returns
\texttt{UNKNOWN}, even if a body condition might otherwise be false. The
distinction prevents an unresolved applicability question from being concealed
inside an apparently definitive in-scope result.

The current runtime also has a conservative conflict policy. It checks the
static dependency closure of the requested rule before evaluating branches.
Conflicting referenced evidence produces \texttt{CONFLICT}, including when a
later scope branch would otherwise make the body irrelevant. This policy is
part of the declared engineering semantics. It should not be attributed to the
wording of the SFC circular. Chapter~\ref{ch:proof} explains why such a policy
must be fixed before comparing implementations.
% END SOURCE UNIT M02:02

% BEGIN SOURCE UNIT M02:03
\section{The named cases are an argument about the rule}
\label{merge:M02:03}

\label{sec:circular-cases}
The 22 named cases are most useful when read as a sequence of challenges to an
interpretation. They do not merely fill rows in a test fixture. Each changes a
fact or evidentiary condition that a plausible implementation might mishandle.
The baseline case is a component classified as a gift, linked to one specific
fund, with no fee-discount exception, and known in-scope actor and activity.

\subsection{The alternative promotional route}
Case g01 establishes the baseline prohibition trigger. Under
equation~\eqref{eq:gift-body}, $G$ and $P$ are true and $D$ is false, so the
body is true. This is a necessary example but a weak test on its own. An
incorrect expression $G\land P$ would also return true, even though it omits
both an exception and an alternative route.

Case g02 changes the promotional facts: no individual fund is named, but the
component is linked to a fund category. Now $P$ is false and $T$ is true.
The disjunction remains true. A candidate that omitted the product-type route
would return false, making this case a concrete witness to the omission. The
source's inclusion of that route supplies the interpretive objection; the
different program outputs make its consequence inspectable.

Case g15 uses the fund-house example through a supplied reviewed product-type
classification. Its purpose is not to write a text classifier for fund-house
mentions. It checks that the chosen formal vocabulary has a place to represent
the classification when review establishes it. This distinction prevents a
source example from being silently generalized into a universal inference from
an entity name.

Case g09 establishes a specific-product link while leaving the product-type
link unknown. The known route is sufficient for the disjunction. Unknown
information about a second route does not erase the first. Case g14 reverses
the situation: the specific-product route is false and the remaining route is
unknown. Here the unresolved fact can change the outcome, so the correct result
under the declared semantics is unknown.

Taken together, these cases explain the meaning of ``either'' in an executable
rule with partial evidence. The system must neither require every route to be
known true nor convert every missing route to false. This is a semantic
distinction, not a preference about how verbose the explanation should be.

\subsection{The exception and the object it qualifies}
Case g03 changes the baseline component to one classified as solely a fee or
charge discount. The exception settles the selected body condition as false.
The result says that this particular gift restriction is not triggered under
the supplied classification. It does not resolve advertising, authorization or
any other requirement associated with the offer.

Case g04 returns to a voucher while another component of the same offer is a
fee waiver. The assessed component remains the voucher, so the waiver does not
set its $D$ value to true. This case challenges an implementation that stores a
single offer-level discount flag and reuses it for every component. If that flag
is the only available data, the adapter has insufficient information to apply
the proposed component-level rule. It should request a finer assessment rather
than derive a convenient value.

Case g07 leaves the discount classification unresolved for a component otherwise
meeting the trigger. If $D$ is false, the prohibition is triggered; if $D$ is
true, the exception applies. Both possibilities remain consistent with the
supplied evidence. Unknown is therefore the informative answer. A system that
turns this unknown into false would remove the exception by default; one that
turns it into true would invent the exception.

Case g16 provides known discount evidence while leaving the gift and promotional
classifications unknown. Within known scope, a true $D$ makes $\neg D$ false,
which settles this conjunction as false. The case demonstrates that every
missing fact need not prevent every narrow conclusion. Its safety depends on
the specific conclusion being reported. It would be unsafe to broaden the
result into a complete clearance of the marketing arrangement.

The contrast between g04 and g16 is particularly useful for reviewer training.
An established exception can settle a condition, but only for the object to
which it applies. A precise value attached to the wrong subject is not better
evidence than an unknown value attached to the right one.

\subsection{Classification is not supplied by the formula}
Case g08 concerns additional returns whose gift classification has not been
settled. The circular's qualified treatment of such returns motivates an
assessment; it does not establish $G$ for every arrangement. The proposed
candidate must retain unknown when the remaining facts would make that
classification decisive. Replacing the missing judgment by an unconditional
``additional return implies gift'' would create a new interpretive rule.

Cases g06 and g22 establish that no gift is offered. In g06 the discount
classification is also missing. The false gift predicate settles the body as
false under the declared Boolean semantics. Case g22 makes the point with a
known product link: a promotional link alone does not establish that a gift
exists. These cases catch a candidate that assumed $G$ from the context instead
of requiring evidence for it.

Case g05 establishes that neither product-link route holds. This also makes the
body false, while leaving other regulatory questions outside the result. It
does not establish that every incentive unrelated to an identified product is
permissible under all applicable requirements. The selected condition is only
one element of a broader assessment.

Case g17 separates regulatory classifications. A separate assessment finds no
structured product, but the component is a non-discount gift with a relevant
product link. The gift trigger remains true. The case challenges a workflow
that exits early after resolving one concern and never evaluates another.
Correct integration therefore requires an inventory of outstanding controls,
not a single reassuring flag named \texttt{regulatory\_check\_passed}.

\subsection{Applicability precedes an in-scope answer}
Cases g10 and g11 establish that, respectively, the actor and the product/activity
are outside the selected profile. The response is out of scope. The underlying
Code requirement may have a broader reach; the example cannot answer that
broader question. Treating out of scope as false would erase the distinction
between an in-scope condition that is not triggered and a question the selected
control does not purport to decide.

Case g12 leaves the fund/activity scope unknown. The system cannot establish
whether the profile applies. Case g13 leaves actor scope unknown while
establishing no gift. A Boolean programmer might be tempted to use the false
gift fact to return false immediately. The declared profile instead reports
unknown scope. This preserves the fact that the system has not established the
context in which its body answer would be meaningful.

This is a design choice that an independent implementation must reproduce.
Another language could define a different scope-combination policy, but it
would then implement a different specification. Differential testing only has
a clear target after the policy has been made explicit. A reviewer must also
decide how an unknown scope result affects the transaction workflow.

\subsection{Contradiction, expiry and later knowledge}
Case g18 supplies conflicting gift evidence. The result is conflict, not
unknown and not the value favored by the latest-arriving record. Unknown means
the evidence does not determine a value. Conflict means admissible records
assert incompatible values. These require different remediation: obtaining
missing evidence in the first case, and resolving or superseding contradictory
evidence in the second.

Case g19 combines an out-of-scope actor with conflicting discount evidence.
The current static-closure conflict policy returns conflict before scope
evaluation. This is deliberately conservative. The example must teach the
policy honestly because a reviewer might otherwise expect out of scope and
mistake the difference for a Java bug. The policy is recorded as engineering
semantics, subject to separate operational review.

Case g20 makes the gift evidence stale at assessment time. The rule does not
continue treating an expired classification as current. Case g21 places the
recording time after the historical knowledge cutoff. A historical assessment
cannot use that later knowledge. Both become unknown when the missing gift
classification remains decisive, but their reason codes and repair actions
differ. Refreshing an expired record does not justify rewriting what was known
at an earlier point in time.

All 22 cases therefore contribute a distinct aspect of the specified behavior.
Some concern the wording of the selected restriction. Others concern evidence,
scope and operational conventions needed to execute the reading. The case
packet should label those sources of expected behavior separately. Otherwise
an implementation convention can acquire the apparent authority of a legal
sentence simply because both are stored in the same JSON file.
% END SOURCE UNIT M02:03

% BEGIN SOURCE UNIT M02:04
\section{What the 729 combinations cover}
\label{merge:M02:04}

\label{sec:circular-enumeration}
Each of the six abstract inputs takes one of three values: true, false or
unknown. The number of assignments is therefore
\begin{equation}
 3^6=729.
 \label{eq:case-domain-size}
\end{equation}
This is a statement about a finite abstraction. It does not count legal
interpretations, marketing documents, evidence histories or real client offers.
The enumeration is exhaustive for that abstraction because it visits every
combination of the six declared values.

The independent truth oracle handles scope first. For the four body facts it
replaces each unknown by both possible Boolean values and evaluates the formula
under every completion. If every completion returns true, the oracle returns
true; if every completion returns false, it returns false; otherwise it returns
unknown. The code does not import RuleIR or the reference evaluator. This makes
it useful for detecting a defect in how the expression was encoded or evaluated.

For this particular body, each input occurs once and with a single polarity.
The completion calculation agrees with the current three-valued operations.
That agreement should not be generalized to every expression. For example, a
formula containing a variable and its negation can be true under every Boolean
completion while a compositional three-valued evaluator returns unknown. The
next chapter derives that difference. The oracle's representation is independent
of the candidate AST, but its semantic scope remains specific.

Conflict, stale evidence and later-recorded evidence are not extra members of
the three-value enumeration. Their normalization and precedence are exercised
by named cases and other application checks. The count 729 must not be used to
imply exhaustive testing of those mechanisms. Nor does it exhaust numerical
inputs, dates, arbitrary rule graphs or the Java runtime's entire implementation.
A result summary should always name the finite object that was exhausted.

The comparison checked full semantic responses and relevant traces, while
respecting each backend's own identity. It would be wrong to require the Python
and Java engine identifiers to be identical: those identifiers describe
different implementations. It would also be wrong to ignore all metadata in a
comparison, because source, evidence and rule identities are part of an
explainable decision. Chapter~\ref{ch:verification} specifies the distinction
between semantic agreement and implementation-specific provenance.
% END SOURCE UNIT M02:04

% BEGIN SOURCE UNIT M02:05
\section{Mutations ask whether the tests can notice an error}
\label{merge:M02:05}

\label{sec:circular-mutations}
Successful cases can be uninformative if they would also pass a wrong program.
Mutation testing addresses this by deliberately changing the rule and asking
whether the checks detect the change. A mutation is not a plausible legal
interpretation merely because it can be compiled. It is an instrument for
testing the sensitivity of the evidence.

The first retained mutation removed the fee-discount exception. Cases g03, g07,
g16 and g19 detected the change. Some of these detections concern the logical
exception, while g19 also exposes the effect of changing the static dependency
closure. Removing an expression can remove a referenced fact from conflict
analysis. A mutation report should therefore explain why a case detects a
change rather than treating every detection as the same kind of evidence.

The second mutation omitted the product-type route. Cases g02, g14 and g15
detected it. Their relationship to the source is direct: they isolate the
route that the mutation removed, including a case where that route is
unresolved. A baseline case with a named fund could not provide this evidence.

The third mutation assumed that a gift existed without requiring its evidence.
Cases g06, g08, g18, g20, g21 and g22 detected the resulting change. The set
includes known absence, missing classification, conflict, expiry and later
knowledge. This shows why evidence-handling cases are necessary even for a
small logical condition. A wrong classification default can alter outcomes
across several superficially different scenarios.

The fourth mutation ignored the selected actor/product scope. Cases g10 through
g13 detected it. The expected difference includes both out-of-scope and unknown
results. If the test harness compared only a final allow/deny bit after
flattening all other statuses, it could conceal these errors.

All four mutations were compiled into Java before comparison. The result thus
concerns executable altered candidates, not a textual inspection of a proposed
patch. Still, detecting four mutations does not establish detection of every
possible implementation or interpretation defect. Mutation operators are chosen
by people or software and can themselves omit an important failure mode. The
ensemble evaluation must maintain a documented fault model and extend it when
new failures are discovered.
% END SOURCE UNIT M02:05

% BEGIN SOURCE UNIT M02:06
\section{A new discrepancy: what counts as a discount?}
\label{merge:M02:06}

\label{sec:circular-cashback}
The existing tests assume that a reviewer can supply the discount classification.
To investigate the missing interpretation step, consider a synthetic offer in
which a customer pays a fee and later receives cash back. Two candidate
definitions might be proposed. One restricts the exception to a reduction of
the fee initially charged. Another includes a documented refund of a fee that
was actually paid. The source packet alone has not established that both
readings are legally tenable. Their purpose here is to expose a question whose
answer affects implementation.

Fix the actor, activity and promotional link, and assume for the comparison
that the component otherwise meets the gift condition. Let the fee paid be
HK\$100 and the later rebate also HK\$100. Under the first proposed definition,
the later transfer is outside the defined exception. Under the second, it may
fall within the exception if the required refund relationship is established.
Those conditions must be derived from authority or review; the example does
not introduce them as new SFC requirements.

The next action should be a targeted investigation of the applicable definition,
Code and FAQ, not a request for the same model to vote again on the word
``discount.'' A source that explicitly includes the relevant refund treatment
would support refinement of the second candidate. A source that excludes it
would defeat that candidate in the applicable scope. A source discussing
another product, actor or period may be informative without settling the
question. The retrieval result needs an applicability assessment before it can
resolve the discrepancy.

Now alter the cash amount to HK\$150. If a candidate definition treats only a
refund of an actual HK\$100 fee as falling within the exception, the extra
HK\$50 raises a further classification question. It does not follow that the
system can automatically split the transfer into a lawful and unlawful part.
The legal unit of assessment and treatment of mixed benefits must first be
settled. An arithmetic decomposition is easy; its legal significance is the
question under investigation.

A second variation changes the evidence. The marketing material describes the
transfer as a fee refund, but no record establishes that a fee was paid. The
candidate's definition may be clear while its decisive factual premise remains
unknown. This discrepancy belongs to fact classification or evidence
availability, not necessarily to interpretation. The resolution controller
should request the relevant evidence or return unresolved, rather than
changing the definition to make the missing record unnecessary.

A third variation combines a refund with a separate voucher. A whole-offer
exception would remove the need to assess the voucher independently. The
proposed component-level interpretation would retain it. This is a different
branch in the hypothesis space from the timing of a refund. Preserving separate
issue identifiers prevents a source that resolves the timing question from
being incorrectly treated as resolving the mixed-offer question as well.

These variations show how an ensemble can make review more precise. It should
produce a small set of materially different definitions, the source reasons
for each, and the smallest scenarios on which they diverge. The reviewer then
decides an identifiable question. A long narrative saying that an offer appears
reasonable would provide less actionable evidence.
% END SOURCE UNIT M02:06

% BEGIN SOURCE UNIT M02:07
\section{A complete dry run of the proposed workflow}
\label{merge:M02:07}

\label{sec:circular-dryrun}
The following dry run specifies proposed behavior. It is not a report that the
new ensemble has already run. Start by creating an assessment whose context
contains the retained circular revision, declared distributor/fund scope,
selected effective interval and an inventory of dependencies. The historical
engineering example uses a synthetic September 2026 interval; the workflow must
not infer a regulatory commencement date from that setting.

The source-review path inventories the circular and flags the incorporated Code
and FAQ as unresolved. The interpretation path proposes the component-level
gift rule. A second interpretation path omits the product-type route. A
definition-focused path asks whether cashback is a discount and whether a
combined offer changes the assessed unit. At this stage the system has a
candidate set and several issues, not a winner.

The comparator identifies a decision difference between the complete promotional
condition and the candidate omitting the product-type route. It constructs a
case with a gift linked only to a fund category and no exception. The
source-grounding check points to the relevant paragraph. The candidate author
can propose a repair adding the missing route. The system stores the repaired
candidate as a new revision, reruns the discrepancy case and the retained
regression cases, and preserves the original candidate and its rejection reason.

The cashback definition issue follows a different route. Retrieval attempts to
obtain the incorporated sources with the right identity and version. If the
sources are unavailable, another language-model explanation cannot close the
issue. The controller records the failed acquisition and the remaining
dependency. It may retry an authorized retrieval route within its budget, or
stop automatic work because no permitted action can supply the missing
authority. The issue remains unresolved either way.

The mixed-offer issue can produce a useful witness even before legal resolution.
The comparison holds the voucher facts fixed and toggles the presence of a
separate fee waiver. If a candidate's voucher decision changes solely because
of the other component's exception, the report explains the dependence. A
reviewer must then confirm the unit to which the exception applies. The witness
does not grant the controller authority to make that legal decision itself.

The accepted formal candidate can be compiled and checked in parallel with this
research. Candidate compilation is valuable because it exposes implementation
defects early. It does not change the issue status. The report may legitimately
say that the candidate passed all specified execution checks and remains
unavailable for release because source interpretation is unresolved. Keeping
those two observations together is more informative than suppressing either.

When automatic work reaches its finite limit, the controller freezes a report.
It identifies the source revision, surviving definitions, rejected encodings,
discriminating scenarios, completed actions, failed actions, unresolved
dependencies and the exact effect of each unresolved issue on the control.
There is no final vote to select a legally preferred answer. The state is ready
for targeted human review or blocked on a named dependency, according to the
reason for stopping.

A designated reviewer can later resolve an issue using newly obtained authority.
That creates a new context revision. Dependent candidate scores and comparisons
are invalidated where appropriate, and the formal rule and tests are updated.
The earlier report remains a faithful account of what was unresolved before the
new material arrived. This is essential for reproducibility: a current reviewer
must be able to see why the earlier system abstained.
% END SOURCE UNIT M02:07

% BEGIN SOURCE UNIT M02:08
\section{From the JSON oracle to independent adjudication}
\label{merge:M02:08}

\label{sec:circular-independent}
The filename \texttt{oracle.json} is an engineering term. It designates a source
of expected results for tests. It does not establish institutional authority or
legal independence. The file states that the same agent authored it and the
candidate. A future user interface should display that authorship rather than
letting the label ``oracle'' imply a gold standard.

Independent adjudication should begin from the retained source packet and
declared task, before exposing reviewers to the proposed formal expression.
One reviewer identifies requirements, definitions and difficult cases. Another
can independently develop or challenge the interpretation. Disagreements are
recorded and, where possible, resolved with source reasons. A third role may
adjudicate material disagreement. Independence here concerns information and
responsibility as well as the identity of a person or model.

After that source-based stage, reviewers can inspect the candidate rules and
their discriminating cases. This order reduces anchoring on the generated
answer. It does not make reviewers infallible. The reference packet must allow
an unresolved category, later correction and competing defensible readings.
If a reviewer cannot determine the treatment of a cash rebate without the
referenced FAQ, the correct reference outcome is not a guessed Boolean chosen
to make benchmark scoring convenient.

The original 22 cases should remain frozen as historical engineering evidence.
New adjudicated cases belong in a separately versioned packet with its own
authors, source context and review decisions. A changed expected outcome should
explain whether the earlier interpretation was rejected, the scope changed,
new evidence became available or the evaluation semantics were corrected.
Those causes imply different repairs and different consequences for existing
releases.

The chapter has now followed a source to a rule, tested the rule, exposed
unresolved definitions and specified how a future ensemble should handle them.
The remaining task is to state precisely what each formal check can prove.
Without that precision, a solver's success could once again be mistaken for
the missing judgment about meaning.
% END SOURCE UNIT M02:08

